\documentclass[a4paper,12pt]{article}
\usepackage[utf8]{inputenc}
\usepackage{amsmath, amssymb}
\usepackage{graphicx}
\usepackage{geometry}
\usepackage{enumitem}
\usepackage{booktabs}
\geometry{margin=1.5cm}
\usepackage{hyperref}
\setlength{\parindent}{0pt}
\setlength{\parskip}{0.8em}

\title{Autonomous Systems on Mars}
\author{Elena-Antonia Stoica}
\date{}

\begin{document}
\maketitle

\begin{abstract}
Documentation of an autonomous aircraft's protocol during a dust storm on Mars.
The document covers a full problem analysis, including research findings, assumptions,
and key questions, followed by a proposed solution with decision logic and implementation details.
\end{abstract}

\tableofcontents
\newpage

\section{Introduction}
As part of the engineering team designing the future autonomous aircraft operating on Mars,
I have written a document detailing how to handle the following specific scenario:
\vspace{-0.8em}
\begin{quote}
\textit{The aircraft is mid-flight, a dust storm is approaching and the communication with Earth
has a delay greater than the time of the storm's arrival, thus any attempt to deliver or receive messages
is out of the equation. The aircraft must decide what protocol to follow.}
\end{quote}
\vspace{-0.8em}
In this documentation, I will analyse the problems dust storms present, the risks they introduce,
and what has made exploring Mars difficult in the past. I will present a solution for the scenario
described above, detailing my thinking process.

The documentation is divided into two main parts: Problem Analysis and Proposed Solution.

\section{Problem Analysis}
This section contains my approach to the problem, my assumptions, and my thinking process.
I will also present the difficulties I encountered.

\subsection{First Thoughts}
My first thought after reading the scenario was straightforward:
\vspace{-0.8em}
\begin{quote}
\textit{What is a dust storm on Mars actually like? What damages are there?}
\end{quote}
\vspace{-0.8em}
Before proposing any solution, I wanted to understand the environment. I read about Martian
dust storms \cite{nasa_dust_storms, esa_dust_devils} and found the following:

\begin{itemize}[itemsep=0.3em]
    \item Winds reach up to 160 km/h. However, Mars's atmosphere is so thin that the dynamic
    pressure is significantly lower than an equivalent wind speed on Earth, meaning structural
    damage is less of a concern than it first appears. The wind does, however, present real
    challenges: navigation accuracy degrades, visibility drops significantly due to suspended
    dust, and maintaining a stable flight path requires continuous correction.

    \item Dust particles are slightly electrostatic. This initially raised the question of whether
    they could serve as an alternative energy source, which I later reframed as a potential
    navigation and sensing aid rather than a power solution.

    \item Storms can last anywhere from a few hours to several weeks, making duration
    completely unpredictable. This turned out to be one of the central difficulties of the problem.

    \item Approximately once every three Martian years, a global dust storm covers the entire
    planet. This has significant implications for mission planning.

    \item Storms occur most frequently during Martian summer in the southern hemisphere,
    which suggests routing and scheduling as a first line of prevention.

    \item The primary cause is solar heating of the surface, which lifts fine particles into
    the atmosphere and blocks sunlight — directly threatening solar-powered systems.
\end{itemize}

My immediate instinct after reading was also to ask: \textit{can the aircraft simply land before
the storm arrives?} With 12 minutes of warning, this depends on the aircraft's current altitude,
speed, terrain below, and proximity to a safe landing zone — none of which are specified. This
question shaped my thinking early on, and suggested that the protocol must account for scenarios
where immediate landing is not possible.

This led to a natural first outline of the protocol: detour around the storm if possible,
land before it arrives if not, and if neither is an option, ride it out. A detour is preferable
to landing because it keeps the mission on schedule and avoids the logistics complications of
a grounded aircraft. The rest of my analysis was spent stress-testing each of those three options.

\subsection{Research Findings}
The most important case I found during research was the Opportunity rover, which went silent
in 2018 after a global dust storm blocked its solar panels for months \cite{nasa_opportunity}. Opportunity was not
damaged by wind or particles — it simply ran out of power. This reframed the problem
entirely: the primary threat is not structural damage, it is energy starvation.

This led me to look at what power sources are viable independent of sunlight. The most
realistic option, with proven Mars heritage, is the \textbf{RTG} (Radioisotope Thermoelectric
Generator) — used on both Curiosity and Perseverance rovers \cite{nasa_mmrtg}. An RTG produces continuous
low-level electrical power from radioactive decay, requires no sunlight, has no moving parts,
and cannot be switched off. It is not sufficient for full propulsion, but it can keep avionics,
navigation computers, and sensors alive indefinitely. This makes it the ideal baseline power
source for storm survival.

Regarding storm structure, Martian dust storms do not behave like Earth hurricanes. There is
no calm eye at the center. Turbulence is chaotic and unpredictable throughout the storm.
Winds are generally slower near the surface due to friction with the ground, which makes
lower altitude the most viable option for finding a landing window — not navigating toward
a center.

Flying above the storm is also not a viable option. Martian atmosphere thins too rapidly with
altitude, and most aircraft designs cannot generate sufficient lift beyond a certain ceiling,
which falls well below storm height in major events. Ingenuity, the first autonomous aircraft
to fly on Mars, operated at the equivalent of approximately 30 km altitude on Earth —
already at the edge of what rotorcraft can achieve \cite{nasa_ingenuity}. There is no meaningful
altitude headroom above that.

Finally, radio signals function through dust. This is relevant for communication with nearby
Martian bases, even mid-storm.

\subsection{Key Questions}
The following questions emerged during my analysis as the most important ones to answer
before proposing any solution:

\begin{enumerate}[itemsep=0.3em]
    \item Can the aircraft detour around the storm? How large is the storm?
    \item If a detour is not possible, can it land before the storm arrives?
    \item If neither, how does the aircraft safely enter and navigate within the storm?
    \item Where are winds weakest within the storm?
    \item How long will the storm last?
    \item Is there critical cargo aboard, and how time-sensitive is the delivery?
    \item How far is the destination, and is there a nearby Martian base?
    \item Can the aircraft maintain positional accuracy inside the storm using inertial navigation
    cross-referenced against terrain maps?
    \item Can signals reach a nearby base mid-storm?
\end{enumerate}

\subsection{Assumptions}
The following assumptions were made in the absence of specific mission parameters:

\begin{enumerate}[itemsep=0.3em]
    \item \textbf{Communication with Earth is not relied upon.} The 14-minute round-trip delay
    exceeds the 12-minute warning window. By the time Earth receives the alert and responds,
    the aircraft will already be inside the storm. All decisions are made autonomously.

    \item \textbf{The aircraft is unmanned.} There is no crew aboard. The priority hierarchy is:
    preserve the aircraft, preserve the cargo, resume the mission.

    \item \textbf{The aircraft is structurally built for Martian conditions.} No fragile external
    hardware is assumed. The aircraft can withstand the dynamic pressure of Martian winds.

    \item \textbf{Pre-loaded high-resolution terrain maps are onboard.} The aircraft does not
    need to see the terrain to navigate — it cross-references its inertial position against
    known terrain data.

    \item \textbf{An RTG is standard equipment.} The aircraft carries a radioisotope
    thermoelectric generator as its baseline power source, independent of solar conditions.

    \item \textbf{Radio signals function through dust.} Communication with nearby bases
    is possible even mid-storm.

    \item \textbf{Martian bases and rovers exist within operational range on some missions.}
    This is not guaranteed, and the protocol accounts for cases where no base is nearby.

    \item \textbf{Storm duration cannot be estimated reliably at decision time.} Local and regional
    storm duration is entirely unpredictable. Global storms occur approximately once every three
    Martian years and their windows are known in advance, making them preventable through mission
    scheduling — but once started, even global storms have unpredictable duration. The protocol
    does not assume the storm will clear within any given timeframe.

    \item \textbf{Martian dust storms have no calm eye.} The aircraft cannot navigate toward
    a calmer center as one might with a hurricane.
\end{enumerate}

\subsection{What Makes This Problem Hard}
Several factors combine to make this a genuinely difficult autonomous decision problem:

\begin{itemize}[itemsep=0.3em]
    \item \textbf{Energy starvation, not structural damage, is the primary threat.} Solar panels
    are blocked. The aircraft must operate on RTG power alone, which is sufficient for avionics
    but likely insufficient for full propulsion against strong winds.

    \item \textbf{Storm duration is unpredictable at decision time.} Local and regional storms
    give no reliable forecast. Global storms are predictable in their occurrence window —
    approximately once every three Martian years — but not in their duration once started.
    Any protocol that depends on knowing when the storm will clear is unreliable.

    \item \textbf{Sensor degradation inside the storm is severe.} Optical sensors fail in dust.
    Sonar-based sensors are unreliable due to the thin atmosphere and particle interference.
    The aircraft is effectively navigating blind, relying on inertial systems and terrain maps.

    \item \textbf{Inertial navigation accumulates error over time.} Wind drift causes positional
    uncertainty to grow the longer the aircraft is inside the storm. This creates time pressure
    on finding a safe landing window.

    \item \textbf{There is no calm center to navigate toward.} Unlike a hurricane, Martian
    dust storms are chaotic throughout. There is no guaranteed low-turbulence zone.

    \item \textbf{Flying above the storm is not feasible.} The atmosphere thins too quickly
    with altitude. The aircraft cannot climb above the storm without losing lift entirely.

    \item \textbf{Cargo delivery introduces compounding unknowns.} Whether to dispatch rovers,
    transfer cargo to a nearby base, or hold and await instructions depends on storm duration,
    base proximity, and destination distance — most of which are unknown at decision time.

    \item \textbf{The decision logic must be unambiguous.} The aircraft is autonomous. Every
    branch of the protocol must have clear, codified trigger conditions. There is no room for
    interpretation.
\end{itemize}

\subsection{How My Thinking Evolved}
My first instinct was not to find a solution — it was to understand the environment. Before
anything else, I wanted to know what a Martian dust storm actually does: what it damages,
how long it lasts, how strong the winds are, and what has gone wrong in past missions because
of one.

That research shifted my focus immediately. The wind, despite reaching 160 km/h, exerts far
less force than on Earth due to the thin atmosphere — structural damage was less of a concern
than I initially assumed. The real danger, as the Opportunity rover made clear, was energy
starvation. The storm does not destroy the aircraft. It blinds the solar panels, and without
power, everything else fails.

This reframing changed the entire direction of my thinking. The question was no longer
\textit{how do we survive the storm} but \textit{how do we keep the aircraft powered and
navigating through it}. The RTG answered the power question. With continuous power assumed,
the storm became a navigation and logistics problem rather than a survival one — and that
opened up the full decision tree: detour, land, or ride it out.

Each of those options introduced new questions — about storm size, terrain, cargo, base
proximity, and signal availability — which led to the final insight: \textit{the aircraft is
not the only asset in the system.} A solution that treats the aircraft in isolation misses
the broader infrastructure. When the aircraft cannot complete the mission alone, the question
becomes what other systems can assist, and under what conditions that assistance is worth invoking.

\section{Proposed Solution}

\subsection{Core Thesis}
The aircraft's survival strategy is not to fight the storm or find a replacement power source,
but to avoid it where possible and outlast it where not — using RTG power to keep systems alive,
terrain maps to navigate blind, and the broader Martian infrastructure to complete the mission
when the aircraft alone cannot.

The key insight that shaped this solution is that the storm is not inherently dangerous to a
well-built aircraft. The wind exerts far less dynamic pressure than an equivalent speed on Earth,
and the RTG provides continuous power independent of solar conditions. The real threats are
navigational: degraded sensors, positional drift, and an unknowable storm duration. The protocol
is therefore designed around information management and decision thresholds, not raw survival.

The protocol follows a strict priority order:
\begin{enumerate}[itemsep=0.3em]
    \item \textbf{Detour} — if the storm is mappable and energy reserves are sufficient, fly around it
    \item \textbf{Land} — if detour is not possible but a safe landing zone is reachable before the storm arrives
    \item \textbf{Ride it out} — if neither is possible, minimize power consumption, drift with the storm,
    and attempt landing when conditions allow
\end{enumerate}

Each option has explicit trigger conditions and thresholds, detailed in Section 3.3.

\subsection{Energy Architecture}
The aircraft carries an RTG (Radioisotope Thermoelectric Generator) as its baseline power source.
This is a deliberate design choice informed by Mars mission heritage — both Curiosity and Perseverance
rely on RTGs for exactly this reason: they produce continuous, stable electrical power from radioactive
decay, require no sunlight, have no moving parts, and cannot be switched off. Ingenuity, NASA's
autonomous helicopter on Mars, demonstrated that autonomous rotorcraft flight in the Martian
atmosphere is achievable \cite{nasa_ingenuity} — but it relied on solar power, which dust storms
render unavailable. The RTG removes this dependency entirely. During a dust storm,
when solar panels are completely blocked, the RTG is the only power source that remains fully
operational without any intervention.

The RTG's output is modest — sufficient to power avionics, navigation computers, communication
systems, and sensors, but not sufficient for full propulsion thrust. This is a known constraint
and the protocol accounts for it: the aircraft does not attempt to fight strong winds on RTG power
alone. Instead, it uses passive aerodynamics where possible, reducing powered flight to the minimum
needed for directional control.

The following table summarises the energy sources evaluated and their viability mid-storm:

\begin{table}[h]
\centering
\begin{tabular}{@{}llll@{}}
\toprule
\textbf{Source} & \textbf{Availability} & \textbf{Power level} & \textbf{Feasibility} \\
\midrule
Solar panels        & No         & High (normally)  & Useless in storm \\
RTG                 & Yes        & Low, steady      & Powers avionics and navigation \\
Electrostatic harvesting & Yes   & Very low         & Sensors only, supplementary \\
Wind / passive gliding   & Yes   & Structural, not electrical & Extends flight time \\
Fission reactor     & Delayed    & High             & Not viable for this design \\
\bottomrule
\end{tabular}
\caption{Energy source evaluation during a Martian dust storm}
\end{table}

The realistic picture is a layered power mode: the RTG keeps all critical systems alive,
electrostatic harvesting from dust particles provides a negligible but non-zero supplementary
trickle to sensors, and the airframe uses passive aerodynamics to reduce the propulsion load.
Together, these allow the aircraft to remain operational indefinitely during a storm — not
comfortably, but reliably.

\subsubsection{Storm Boundaries Detection}

Storm boundary detection also relies on existing onboard sensors rather than dedicated storm
mapping hardware. The aircraft uses a combination of the following to estimate storm size
and position in real time:

\begin{itemize}[itemsep=0.3em]
    \item \textbf{Barometric pressure sensors} — pressure gradients shift measurably at storm edges
    \item \textbf{Temperature sensors} — dust storms cause detectable surface and air temperature changes
    \item \textbf{Wind speed and direction sensors} — storm boundary is detectable as wind patterns change sharply
    \item \textbf{Optical sensors} — the visible dust wall at the storm edge is detectable before
    visibility fully degrades
    \item \textbf{Electrostatic sensors} — the storm's electrostatic field creates measurable gradients
    at its boundaries
\end{itemize}

These inputs are fused to produce a storm boundary estimate with an associated confidence value.
This confidence value is a direct input to the detour decision threshold in Section 3.3.

\subsection{Decision Logic}
The aircraft has 12 minutes from storm detection to make and execute a decision. The protocol
evaluates three options in strict priority order. Each option has explicit trigger conditions
that must all be satisfied before that path is taken. If any condition fails, the protocol
falls through to the next option.

\subsubsection*{Data Flow}
Before any decision is made, data moves through the system in the following sequence:

\begin{enumerate}[itemsep=0.3em]
    \item \textbf{Sensors $\rightarrow$ Fusion Layer.} Raw readings from barometric, temperature,
    wind, optical, and electrostatic sensors are collected and fused into a unified storm
    boundary estimate with an associated confidence value.

    \item \textbf{Fusion Layer $\rightarrow$ State Model.} The fused estimate updates the
    aircraft's internal state: storm position, storm width, time to arrival, current energy
    reserves, inertial position, and positional uncertainty.

    \item \textbf{State Model $\rightarrow$ Decision Engine.} The decision engine reads the
    current state and evaluates each protocol option in priority order against its trigger
    conditions.

    \item \textbf{Decision Engine $\rightarrow$ Action.} The selected action — detour, land,
    or enter survival mode — is passed to the flight control system for execution.

    \item \textbf{Action $\rightarrow$ Transmission.} In parallel with execution, the aircraft
    transmits its current state, decision, and position to Earth and any nearby Martian bases.
\end{enumerate}

During survival mode, steps 2 through 5 repeat every 60 seconds as the aircraft continuously
reassesses whether a landing window has opened.

\subsubsection*{Option 1 — Detour}
A detour is the preferred outcome. It keeps the mission on schedule, avoids all storm-related
risks, and requires no cargo contingency planning. The aircraft attempts a detour if and only
if all of the following conditions are met:

\begin{itemize}[itemsep=0.3em]
    \item \textbf{Storm boundary confidence $\geq$ 70\%.} The aircraft's sensor fusion must
    produce a storm width and position estimate with sufficient confidence. Below this threshold,
    committing to a detour route is too risky — the storm may be larger than estimated.

    \item \textbf{Detour distance is within energy budget with a 20\% reserve.} The estimated
    detour route must be completable on current energy reserves, retaining at least 20\% as a
    safety margin. This margin accounts for wind resistance, course corrections, and the energy
    cost of landing at the destination.

    \item \textbf{Time to storm arrival is sufficient to begin the detour.} With 12 minutes of
    warning, the aircraft must confirm it can clear the storm's edge before the storm reaches
    its current position.
\end{itemize}

If all three conditions are met, the aircraft calculates the shortest detour route around the
storm boundary, updates its flight path, and continues toward the original destination.

\subsubsection*{Option 2 — Land Before the Storm Arrives}
If a detour is not possible, the aircraft evaluates whether it can land safely before the storm
arrives. This option is triggered if all of the following conditions are met:

\begin{itemize}[itemsep=0.3em]
    \item \textbf{A safe landing zone exists within reachable distance.} The terrain map must
    identify a flat, obstacle-free area within the aircraft's current glide and flight range.

    \item \textbf{The landing zone is reachable within the remaining time window.} The aircraft
    must be able to reach and land at the zone before the storm arrives.

    \item \textbf{Wind conditions at the landing zone are within safe landing thresholds.}
    If the storm's leading edge has already reached the candidate landing zone, this option
    is disqualified.
\end{itemize}

If all three conditions are met, the aircraft navigates to the landing zone and executes a
controlled landing. It then transmits its position and status to Earth and to any nearby
Martian bases, and holds for further instructions. Note that Earth's response will arrive
no sooner than 14 minutes after the transmission — the aircraft does not wait for a response
before landing. The transmission is informational, not a request for permission.

\subsubsection*{Option 3 — Ride It Out}
If neither detour nor landing is possible, the aircraft enters storm survival mode. This is
not a passive state — it is an active low-power flight mode with continuous reassessment.

The aircraft takes the following steps:

\begin{itemize}[itemsep=0.3em]
    \item \textbf{Switch to minimum power mode.} Non-essential systems are shut down.
    Only navigation, terrain avoidance, communication, and structural control remain active.
    The RTG provides sufficient power for all of these indefinitely.

    \item \textbf{Descend to lower altitude.} Surface friction reduces wind speed near the ground,
    giving the aircraft more control authority. The aircraft descends to a dynamically calculated
    safe altitude based on terrain map data, with a fixed hard floor of 50 metres above the
    highest terrain feature within a 2 km radius. If positional uncertainty exceeds a safe
    threshold, the fixed floor takes precedence over the dynamic calculation.

    \item \textbf{Adopt a passive gliding profile.} Rather than fighting the wind with full
    propulsion, the aircraft orients itself to extract lift from the wind and drift with the
    storm. Propulsion is used only for directional correction, not to maintain position.

    \item \textbf{Continuously monitor for a landing window.} Every 60 seconds, the aircraft
    reassesses whether a safe landing zone is now reachable. As the storm progresses and
    the aircraft drifts, new terrain becomes available. The moment a landing window is
    confirmed — flat terrain, wind speed within threshold, positional confidence sufficient
    — the aircraft exits drift mode and executes a controlled landing.

    \item \textbf{Transmit position continuously.} The aircraft broadcasts its position and
    status at regular intervals so that nearby bases and Earth can track it throughout the storm.
\end{itemize}

\subsubsection*{Altitude Floor Detail}
The altitude floor during drift is calculated as follows:

\begin{itemize}[itemsep=0.3em]
    \item \textbf{Dynamic floor:} $A_{floor} = \max(\text{terrain height within 2 km radius}) + 50\text{m}$,
    updated continuously from terrain map cross-referenced with inertial position estimate.

    \item \textbf{Fixed hard floor:} 50 metres above the highest terrain feature in the
    pre-loaded map for the current region. This activates automatically when positional
    uncertainty $\sigma > 100\text{m}$, overriding the dynamic calculation.
\end{itemize}

The fixed hard floor is a conservative fallback. It may result in a higher-than-necessary
altitude in flat terrain, but it guarantees terrain clearance when the aircraft can no longer
trust its own position estimate.

\subsection{Terrain Avoidance}
Inside a dust storm, the aircraft is effectively blind. Optical sensors fail in suspended dust,
and sonar-based sensors are unreliable in Mars's thin atmosphere due to poor sound propagation
and particle interference. Terrain avoidance therefore relies entirely on two systems working
in combination: pre-loaded terrain maps and inertial navigation.

\textbf{Primary method — map-based navigation.} Before any mission, the aircraft is loaded with
high-resolution terrain maps of its operational area. These maps include elevation data, known
obstacles, and candidate landing zones. Inside the storm, the aircraft does not need to see the
terrain — it needs to know where it is relative to the map. As long as positional accuracy is
sufficient, the aircraft can cross-reference its inertial position against the map and maintain
safe terrain clearance.

\textbf{Inertial navigation and drift error.} Inertial navigation systems track position by
integrating acceleration and rotation over time from a known starting point. They require no
external signals and function in complete sensor blackout. Their weakness is cumulative error:
small measurement inaccuracies compound over time, and wind drift adds further positional
uncertainty. The longer the aircraft is inside the storm, the less confident its position
estimate becomes. This is why the fixed hard floor fallback exists — when positional uncertainty
$\sigma$ exceeds 100 metres, the aircraft stops trusting its dynamic terrain calculation and
reverts to the most conservative safe altitude for the region.

\textbf{Sonar as last-resort backup.} Despite its limitations in thin atmosphere, a low-frequency
sonar system can serve as a close-range backup for terrain detection at very low altitudes —
specifically during final approach to landing when the aircraft is within 20--30 metres of the
surface. At this range and altitude, the atmosphere is at its densest on Mars, improving sonar
reliability marginally. It is not a primary navigation tool, but it provides a final sanity
check before touchdown.

\textbf{Landing zone selection during drift.} As the aircraft drifts, it continuously scans
its terrain map ahead of its projected path. A candidate landing zone must satisfy the following
conditions before the aircraft commits to it:

\begin{itemize}[itemsep=0.3em]
    \item Terrain slope below a safe threshold within the map data
    \item No known obstacles within the approach corridor
    \item Positional confidence sufficient to trust the map cross-reference
    \item Wind speed at current altitude within controllable range for final descent
\end{itemize}

\subsection{Pseudocode}
The following pseudocode describes the full autonomous decision logic from storm detection
through to resolution. Function names and conditions correspond directly to the thresholds
defined in Sections 3.3 and 3.4.

\begin{verbatim}
# Constants
DETOUR_CONFIDENCE_THRESHOLD = 0.70
ENERGY_SAFETY_MARGIN = 0.20
ALTITUDE_HARD_FLOOR = 50          # metres above highest terrain in region
POSITION_UNCERTAINTY_LIMIT = 100  # metres, triggers hard floor fallback
DRIFT_REASSESSMENT_INTERVAL = 60  # seconds
SONAR_ACTIVATION_ALTITUDE = 30    # metres, last-resort terrain check

# --- Storm Detection ---
function on_storm_detected(storm_data, aircraft_state):
    transmit_status_to_earth(aircraft_state, storm_data)
    transmit_status_to_nearest_base(aircraft_state, storm_data)

    decision = evaluate_protocol(storm_data, aircraft_state)
    execute(decision)


# --- Top-Level Protocol ---
function evaluate_protocol(storm_data, aircraft_state):

    if can_detour(storm_data, aircraft_state):
        return plan_detour(storm_data, aircraft_state)

    if can_land_before_storm(storm_data, aircraft_state):
        return plan_immediate_landing(aircraft_state)

    return enter_survival_mode(aircraft_state)


# --- Option 1: Detour ---
function can_detour(storm_data, aircraft_state):
    confidence = estimate_storm_boundary_confidence(storm_data)
    if confidence < DETOUR_CONFIDENCE_THRESHOLD:
        return False

    detour_route = calculate_detour_route(storm_data, aircraft_state.position)
    detour_cost  = estimate_energy_cost(detour_route, aircraft_state)
    available    = aircraft_state.energy_reserves * (1 - ENERGY_SAFETY_MARGIN)
    if detour_cost > available:
        return False

    time_to_clear = estimate_time_to_clear_storm_edge(detour_route, aircraft_state)
    if time_to_clear > aircraft_state.time_to_storm_arrival:
        return False

    return True

function plan_detour(storm_data, aircraft_state):
    route = calculate_detour_route(storm_data, aircraft_state.position)
    return { action: "DETOUR", route: route }


# --- Option 2: Land Before Storm ---
function can_land_before_storm(storm_data, aircraft_state):
    zone = find_nearest_safe_landing_zone(
               aircraft_state.position,
               aircraft_state.terrain_map)
    if zone is None:
        return False

    time_to_land = estimate_flight_time(aircraft_state.position, zone)
    if time_to_land > aircraft_state.time_to_storm_arrival:
        return False

    if storm_leading_edge_reached(zone, storm_data):
        return False

    return True

function plan_immediate_landing(aircraft_state):
    zone = find_nearest_safe_landing_zone(
               aircraft_state.position,
               aircraft_state.terrain_map)
    return { action: "LAND", target_zone: zone }


# --- Option 3: Survival Mode ---
function enter_survival_mode(aircraft_state):
    shutdown_nonessential_systems()
    activate_passive_glide_profile()
    descend_to_safe_altitude(aircraft_state)

    while True:
        transmit_position(aircraft_state)
        wait(DRIFT_REASSESSMENT_INTERVAL)

        update_position_estimate(aircraft_state)
        update_altitude_floor(aircraft_state)

        zone = find_nearest_safe_landing_zone(
                   aircraft_state.position,
                   aircraft_state.terrain_map)

        if zone is not None and landing_conditions_met(zone, aircraft_state):
            return { action: "LAND", target_zone: zone }


# --- Altitude Management ---
function descend_to_safe_altitude(aircraft_state):
    floor = compute_altitude_floor(aircraft_state)
    aircraft_state.target_altitude = floor

function compute_altitude_floor(aircraft_state):
    if aircraft_state.position_uncertainty > POSITION_UNCERTAINTY_LIMIT:
        # positional confidence too low — use conservative fixed floor
        return get_regional_max_terrain_height() + ALTITUDE_HARD_FLOOR

    # dynamic floor — terrain map cross-referenced with current position
    local_max = max_terrain_height_within_radius(
                    aircraft_state.position, radius=2000)
    return local_max + ALTITUDE_HARD_FLOOR


# --- Landing Execution ---
function execute_landing(target_zone, aircraft_state):
    navigate_to(target_zone)

    if aircraft_state.altitude < SONAR_ACTIVATION_ALTITUDE:
        activate_sonar_terrain_check()   # last-resort sanity check

    perform_controlled_landing(target_zone)
    transmit_status_to_earth(aircraft_state, status="LANDED")
    transmit_status_to_nearest_base(aircraft_state, status="LANDED")
\end{verbatim}

\subsection{Cargo Contingency}
The base protocol described in Section 3.3 applies regardless of cargo. However, if the aircraft
is carrying cargo flagged as mission-critical in its pre-loaded mission parameters, an additional
layer of decision logic applies after landing.

\subsubsection*{Why landing is still the right first step}
Even with critical cargo aboard, the aircraft should not attempt to fly through the storm
indefinitely. Storm duration is unpredictable — it could last weeks. Drifting with critical
cargo for an unknown period, with degraded navigation and no guarantee of reaching the
destination, introduces more risk than a controlled landing. The protocol therefore lands
as early as possible, then evaluates cargo delivery options from the ground.

\subsubsection*{Cargo delivery decision tree}
Once the aircraft has landed, the following logic applies:

\begin{enumerate}[itemsep=0.3em]
    \item \textbf{Transmit position, cargo status, and destination to Earth and nearest base.}
    Radio signals function through dust, so this is reliable regardless of storm conditions.
    Earth's response will arrive in approximately 14 minutes — the aircraft does not wait
    for it before evaluating ground options.

    \item \textbf{Check proximity of nearest Martian base.}
    If a base is within rover range, the aircraft signals the base to dispatch a rover.
    Rovers equipped with RTGs or batteries can operate in dust storms — they are low to the
    ground, do not depend on lift, and are not affected by visibility in the same way an
    aircraft is.

    \item \textbf{Evaluate delivery feasibility based on destination distance.}
    \begin{itemize}[itemsep=0.3em]
        \item If the destination is within approximately one week of rover travel time,
        dispatch a rover to collect and deliver the cargo directly.
        \item If the destination is further, the nearest base retrieves the cargo and holds
        it, then signals the destination to arrange onward transport once the storm clears.
        \item If no base is within range, the aircraft holds the cargo, transmits its
        position continuously, and awaits instructions from Earth or the nearest base.
        This is the fallback of last resort — it introduces delay, but preserves the cargo
        and the aircraft.
    \end{itemize}
\end{enumerate}

\subsubsection*{The storm duration problem}
A key difficulty here is that rover dispatch without knowing storm duration risks wasting
power. A rover sent into a multi-week global storm may not survive the journey. The protocol
therefore includes a storm type check before dispatching rovers:

\begin{itemize}[itemsep=0.3em]
    \item \textbf{If the storm is identified as a local or regional storm} — duration is
    unknown but statistically likely to resolve within days to a few weeks. Rover dispatch
    is evaluated based on destination distance and rover energy budget.
    \item \textbf{If the storm is identified as a global storm} — duration could extend
    to months. Rover dispatch is suspended. The aircraft holds cargo and awaits Earth
    instructions. Global storm windows are known in advance from orbital monitoring, so
    this scenario should have been prevented at the mission planning stage.
\end{itemize}

Storm type is identified from pre-loaded orbital weather data updated before departure,
cross-referenced with real-time sensor readings on storm scale and pressure gradients.

\subsubsection*{Pseudocode extension for cargo}
\begin{verbatim}
# Additional constant
ROVER_MAX_TRAVEL_DAYS = 7

# Called after any landing when cargo flag is set
function handle_cargo_contingency(aircraft_state, cargo):
    if not cargo.is_mission_critical:
        return   # standard hold, await aircraft resumption

    transmit_cargo_status(aircraft_state, cargo)

    storm_type = identify_storm_type(aircraft_state.sensor_data)
    if storm_type == "GLOBAL":
        # do not dispatch rovers, hold and await Earth instructions
        hold_and_await_instructions(aircraft_state, cargo)
        return

    nearest_base = find_nearest_base(aircraft_state.position)
    if nearest_base is None:
        hold_and_await_instructions(aircraft_state, cargo)
        return

    travel_days = estimate_rover_travel_time(
                      nearest_base.position,
                      cargo.destination)

    if travel_days <= ROVER_MAX_TRAVEL_DAYS:
        signal_base_dispatch_rover(nearest_base, aircraft_state, cargo)
    else:
        signal_base_retrieve_and_hold(nearest_base, aircraft_state, cargo)


function hold_and_await_instructions(aircraft_state, cargo):
    while not instructions_received():
        transmit_position(aircraft_state)
        transmit_cargo_status(aircraft_state, cargo)
        wait(DRIFT_REASSESSMENT_INTERVAL)
\end{verbatim}

\subsection{Why This Approach}
The solution presented here was not the first one I considered. My initial instinct was to
understand what the storm could actually do — what damages it causes, how strong it is, how
long it lasts. Only after understanding the environment did the energy problem become central.
The Opportunity rover made it concrete: a machine that was not destroyed by a storm, but simply
ran out of power waiting for one to clear. The real problem was never the storm itself. It was
the assumptions baked into the system — that solar power would always be available, and that
the mission could continue uninterrupted.

This shaped the entire approach. Rather than designing around a single failure mode, the
protocol is designed around uncertainty: unknown storm duration, degraded sensors, unreliable
positioning, and an infrastructure that may or may not be within range. Every decision branch
has an explicit fallback. Every threshold is named and adjustable. Nothing in the protocol
assumes information the aircraft does not have.

There are several specific reasons why this approach is preferable to the alternatives:

\textbf{RTG over solar dependency.} An aircraft designed for Mars that depends solely on solar
power is fragile by design. Dust storms are not rare — they are a predictable feature of the
Martian environment. The RTG removes the single most dangerous dependency in the system.
It does not solve every problem, but it ensures the aircraft can always think, navigate, and
communicate, regardless of conditions.

\textbf{Detour over avoidance.} The protocol prefers a detour to simply waiting on the ground.
A grounded aircraft waiting for a storm to clear — with no estimate of when that will be —
is a mission that may never resume. A detour keeps the mission moving. It is only attempted
when the conditions genuinely support it, but it is always the first option evaluated.

\textbf{Outlasting over fighting.} When neither detour nor landing is possible, the aircraft
does not attempt to maintain its position or push through the storm on propulsion alone. RTG
power is insufficient for sustained thrust against strong winds. Fighting the storm would
drain the aircraft's limited power reserves and leave it in a worse position. Passive gliding
conserves energy, extracts lift from the wind, and buys time — which is exactly what the
situation requires.

\textbf{Systems thinking over aircraft-only thinking.} The cargo contingency section reflects
a broader principle: the aircraft is one node in a larger system. Rovers, bases, and orbital
monitoring all exist. A solution that ignores them is incomplete. The protocol uses them where
appropriate, without over-relying on them — because their availability is not guaranteed.

\textbf{Unambiguous logic over flexible judgment.} Every branch of the protocol has clear,
codified trigger conditions. This is not a stylistic choice — it is a requirement. An
autonomous system cannot exercise judgment. It can only evaluate conditions and execute
instructions. Ambiguity in the protocol is a bug, not a feature. Every threshold in the
pseudocode exists because a vague instruction like ``land if it seems safe'' is not
implementable.

The honest limitation of this approach is that it depends on the quality of the terrain maps
and the accuracy of the inertial navigation system. If the maps are incomplete or the
positional error grows too large, the aircraft's ability to find a safe landing zone degrades.
This is acknowledged in the fixed hard floor fallback, but it cannot be fully mitigated in
software — it is ultimately a hardware and pre-mission data quality problem.

\section{Conclusion}
This document set out to answer a deceptively simple question: what should an autonomous
aircraft do when a dust storm is 12 minutes away and Earth cannot help in time?

The answer required understanding the environment before designing a solution. Martian dust
storms are not structurally dangerous to a well-built aircraft — they are energetically and
navigationally dangerous. The Opportunity rover did not fail because the storm was violent \cite{nasa_opportunity}.
It failed because the system had a single point of failure: solar power. That insight drove
every decision in this protocol.

The proposed solution is built on three principles. First, remove the single point of failure
— an RTG \cite{nasa_mmrtg} provides continuous power independent of solar conditions, ensuring the aircraft can
always think, navigate, and communicate. Second, prefer the option that keeps the mission
moving — detour first, land second, ride it out only when necessary. Third, design for
uncertainty — unknown storm duration, degraded sensors, and variable infrastructure mean
every decision branch needs an explicit fallback, not an assumption that conditions will
be favourable.

The result is a protocol that is not optimistic. It does not assume the storm will be short,
that a landing zone will be nearby, or that a base will be within rover range. It assumes the
worst and plans accordingly. When conditions are better than expected, the protocol naturally
takes the better path. When they are not, it has a fallback.

The honest limitation is that the protocol's reliability is bounded by the quality of the
terrain maps and the accuracy of the inertial navigation system. These are pre-mission
dependencies that cannot be compensated for in software alone. A mission that departs with
incomplete maps or poorly calibrated inertial systems has a weaker protocol, regardless of
how well the decision logic is written.

Avoiding the southern hemisphere in Martian summer, scheduling around known global storm windows,
and maintaining awareness of the three-year cycle are all upstream decisions that cost nothing
and eliminate entire categories of risk. The best autonomous decision is the one the aircraft
never has to make.

\newpage
\bibliographystyle{plain}
\bibliography{References}

\end{document}
